% Converted from a PDF by kakemono. Compiles with pdfLaTeX, XeLaTeX or
% LuaLaTeX; the pictures are files in a folder next to this one.
\documentclass[11pt]{article}
\usepackage{iftex}
\ifPDFTeX
  \usepackage[T1]{fontenc}
  \usepackage{lmodern}
\else
  \usepackage{fontspec}
\fi
\usepackage[margin=2.5cm]{geometry}
\usepackage{amsmath,amssymb,mathtools}
\usepackage{graphicx}
\usepackage{booktabs}
\usepackage{multirow}
\usepackage[colorlinks,allcolors=blue]{hyperref}
% Relations written as MathJax names them.
\providecommand{\lt}{<}
\providecommand{\gt}{>}
% A character that the fonts of pdfLaTeX do not have (XeLaTeX and LuaLaTeX
% write it).
\newcommand{\kakemonomissing}{\textbf{?}}
% A picture of an area of the PDF page: #1 is its width in the PDF, #2 the
% file. It is at most as wide as the text and 80% as high.
\newlength{\kakemonowidth}
\newcommand{\kakemonopicture}[2]{%
  \if\relax\detokenize{#2}\relax\else
    \setlength{\kakemonowidth}{#1}%
    \ifdim\kakemonowidth>\linewidth\setlength{\kakemonowidth}{\linewidth}\fi
    \includegraphics[width=\kakemonowidth,height=0.8\textheight,keepaspectratio]{#2}%
  \fi}
% A table, shrunk to the width of the text when it is wider.
\newsavebox{\kakemonobox}
\newenvironment{kakemonotable}
  {\begin{lrbox}{\kakemonobox}}
  {\end{lrbox}%
   \ifdim\wd\kakemonobox>\linewidth
     \resizebox{\linewidth}{!}{\usebox{\kakemonobox}}%
   \else
     \usebox{\kakemonobox}%
   \fi}
